% SPDX-FileCopyrightText: © 2026 Idiap Research Institute <contact@idiap.ch>
%
% SPDX-FileContributor: Andre Freitas <andre.freitas@idiap.ch>
% Neuro-symbolic AI Group
%
% SPDX-FileContributor: Danilo Gusicuma <danilo.gusicuma@idiap.ch>
% Neuro-symbolic AI Group
%
% SPDX-License-Identifier: MIT
\subsection{Abductive proof work: making a proposed route answerable}
\label{app:granular-abducer-work-loop}
An active proof gap is an obligation, not yet a route for meeting it. The four
occurrences make four different commitments. Route formation ($a$) proposes a
hypothesis that could close the gap. Claim derivation ($d$) takes that hypothesis
as a premise and exposes a consequence. Refutation ($r$) challenges that exact
consequence. Routing ($e$) retains both the consequence and its assessment to
choose what work is justified next. This retained join matters: an adverse test
without the claim cannot say which part of the proposed route failed.

\paragraph{How the search is constrained.}
The linked abductive skill first attempts deduction to locate the exact missing
support, then constructs competing small sets of additional premises. It compares
explanatory coverage, compatibility with accepted knowledge, simplicity, and
novelty risk, and attaches a falsifier to each set. Thus the search space is not
``any plausible explanation'': it is alternative repairs of a named entailment
gap. The linked outer work procedure alternates discovery and proof attempts,
reassesses end-to-end plausibility, and permits persistence only when failure
sharpens a premise, object, regime, or bridge. Repeating the same blocker instead
requires backtracking or a pivot. Its machinery assessment asks which method
actually controls the bottleneck, excluding a prestigious analogy with no
transfer obligation. These are prescribed skill procedures; the four native
occurrences do not implement that persistent controller.

Here $G,H,C,R,S$ alias the native gap, route hypothesis, derived claim,
refutation assessment, and routing assessment types, respectively. Write
$\mathsf{Gen}[X]=\mathtt{Generated}[X]$; $z.\mathtt{value}$ is the declared
projection, not an inferred proof of $X$.
\[
\begin{aligned}
 a &:G\to\mathsf{Gen}[H],&d &:H\to\mathsf{Gen}[C],\\
 r &:C\to\mathsf{Gen}[R],&e &:J\to\mathsf{Gen}[S],\\
 h&=a(g),\quad c=d(h.\mathtt{value}),&v&=r(c.\mathtt{value}),\\
 J&=\{\mathtt{derived\_claim}:C,\ \mathtt{refutation\_assessment}:R\},\\
 s&=e\{\mathtt{derived\_claim}=c.\mathtt{value},\
       \mathtt{refutation\_assessment}=v.\mathtt{value}\}.
\end{aligned}
\]
The relation contracts record assumption, test, and calibration links. In
particular, \emph{Treats\allowbreak Proof\allowbreak Route\allowbreak As\allowbreak Claim\allowbreak Premise} licenses the description
``conditional on this route''; it does not certify that the route holds. The
other two relations attach the challenge and the resulting work assessment to
that commitment. Their source/target records do not supply logical soundness
proofs.

\paragraph{Constructed example.}
To establish a graph-family mixing bound, $a$ proposes a uniform spectral gap.
Under this premise, $d$ derives exponential decay. Testing cycle graphs at $r$
shows their gaps approach zero with size. This refutes the strengthening that
connected graphs have a size-independent gap, not the conditional decay theorem.
The final join supports reframing the
graph family or retargeting the bound, rather than declaring all decay reasoning
false. This is a mathematical illustration, not an observed workflow run.

The native sequence returns a routing assessment with release, reframe,
retarget, or refuted alternatives. The original ERFL additionally switches to
release promotion, conceptual exploration, or target navigation, and stops when
$r$ is refuted. Those transitions, including the early stop before calibration,
are not executable branches in this native sequence.

% SPDX-FileCopyrightText: © 2026 Idiap Research Institute <contact@idiap.ch>
%
% SPDX-FileContributor: Andre Freitas <andre.freitas@idiap.ch>
% Neuro-symbolic AI Group
%
% SPDX-FileContributor: Danilo Gusicuma <danilo.gusicuma@idiap.ch>
% Neuro-symbolic AI Group
%
% SPDX-License-Identifier: MIT
% Generated from step-explanations.json
\paragraph{Step-by-step reconstruction.}
Follow a proposed uniform diffusion bound for du/dt=-Lu over finite connected undirected graphs, using the symmetric combinatorial Laplacian L and mean-zero initial data.
The following is a source-guided illustrative reconstruction. Internal procedures are prescriptions; the native declarations do not establish their execution.
\begin{enumerate}[leftmargin=*,label=\arabic*.]
\item \textbf{Repair the missing premise} (\texttt{route\_\allowbreak{}hypothesis}).
\begin{description}[style=standard,labelindent=0pt,leftmargin=5.5em,labelwidth=5em,labelsep=0.5em,align=left,font=\normalfont\bfseries,itemsep=0.15em,parsep=0pt,topsep=0.3em]
\item[Purpose] Identify a premise sufficient for uniform decay.
\item[Input] Active obligation: decay uniform in graph size
\item[Method] First attempt deduction; locate the absent support, propose small competing premise sets, and attach a falsifier to each.
\item[Check] Compare explanatory coverage and compatibility without treating simplicity as proof.
\item[Output] Route hypothesis: positive common spectral gap
\item[Next] Send this conjectural premise to derivation; incompatible premise sets require revision.
\item[Example] The gap becomes the proposal that every graph in the selected family has first positive eigenvalue at least a common c greater than zero.
\end{description}
\item \textbf{Expose the conditional consequence} (\texttt{derived\_\allowbreak{}claim}).
\begin{description}[style=standard,labelindent=0pt,leftmargin=5.5em,labelwidth=5em,labelsep=0.5em,align=left,font=\normalfont\bfseries,itemsep=0.15em,parsep=0pt,topsep=0.3em]
\item[Purpose] Derive the proposed premise’s conditional payoff.
\item[Input] Proposed common gap; graph-family quantifier
\item[Method] Derive what the added premise would imply while tracking where it is used; do not discharge that premise by using it.
\item[Check] A valid conditional derivation must retain the norm, subspace and common constant.
\item[Output] Exponential-decay claim, conditional on the gap
\item[Next] Refutation receives this claim; a missing spectral lemma remains an obligation.
\item[Example] The proposed gap yields mean-zero norm decay bounded by the initial norm times exp(-ct), conditional on the symmetric Laplacian assumptions.
\end{description}
\item \textbf{Attack the exact strengthening} (\texttt{refutation\_\allowbreak{}assessment}).
\begin{description}[style=standard,labelindent=0pt,leftmargin=5.5em,labelwidth=5em,labelsep=0.5em,align=left,font=\normalfont\bfseries,itemsep=0.15em,parsep=0pt,topsep=0.3em]
\item[Purpose] Test the strengthening without conflating premise and consequence.
\item[Input] Derived decay claim; retained uniform-gap premise
\item[Method] Probe boundary families and seek a counterexample to the implication actually being asserted; distinguish a false premise from an invalid derivation.
\item[Check] A counterexample must satisfy the stated domain and attack the same quantifier.
\item[Output] Refutation finding: connectedness insufficient for a common gap
\item[Next] Pass the targeted finding to calibration; bounded non-refutation supplies no missing proof.
\item[Example] For cycle graphs the first positive eigenvalue is 2-2 cos(2 pi/n), which tends to zero. This defeats a size-independent gap for all connected graphs.
\end{description}
\item \textbf{Choose the next proof task} (\texttt{work\_\allowbreak{}routing}).
\begin{description}[style=standard,labelindent=0pt,leftmargin=5.5em,labelwidth=5em,labelsep=0.5em,align=left,font=\normalfont\bfseries,itemsep=0.15em,parsep=0pt,topsep=0.3em]
\item[Purpose] Choose a sharper task from the localized failure.
\item[Input] Original derived claim + its refutation assessment
\item[Method] Compare the route with the target and identify whether a premise, graph family or bound must change.
\item[Check] Continuation must sharpen the obligation; repeating the same failed premise is not progress.
\item[Output] Routing assessment: bounded reframing or retargeting rationale
\item[Next] Original policy can route onward or stop; the native sequence returns an assessment only.
\item[Example] Retain the conditional decay theorem and replace the unrestricted family claim by the task of finding structural assumptions that imply a uniform gap.
\end{description}
\end{enumerate}
\paragraph{Changed-input contrast.}
If the input instead supplies a proved common spectral-gap bound for a restricted family, the cycle objection no longer defeats that premise. The conditional decay implication can then support work on that family; connectedness alone cannot.


\subsection{Discovery marathon: separating a pattern from its explanation}
\label{app:granular-discovery-marathon-loop}
Discovery here is a bounded comparison of commitments. Scope grounding ($e_0$)
fixes what the run may use. Abduction ($a$) proposes an explanation; induction
($i$) extracts a provisional pattern from that same scope. They are parallel
occurrences because a pattern and an explanatory hypothesis are independently
consumed results. Neither becomes confirmation merely because both were formed.
Model probing ($m$) asks what their combination implies under a model;
contrastive probing ($c$) asks what a discriminating comparison reveals.
Refutation ($r$) combines these different challenges. Calibration ($e_1$)
rejoins the original hypothesis and pattern with the refutation assessment.

\paragraph{What makes a discovery cycle informative.}
The portfolio procedure treats target commitments as fixed constraints and
hypotheses as replaceable. A candidate earns further work by identifying the
commitments it advances and the evidence its next step could change; a novel but
irrelevant branch remains visible without entering the proof queue. The linked
probe procedure asks rivals to predict the same observable under the same
controls before any run. A minimal perturbation of an assumption or carrier then
reveals which dependency the proposed mechanism needs. If rivals cannot be
distinguished, the correct output is non-identifiability, not preference for the
more elaborate model. Negative and boundary outcomes feed branch revision;
repeated failure without new learning forces a change of route. These prescribed
procedures explain the two probe roles, but their current entry gates are richer
than the native hypothesis--pattern join and are not established by its types.

With $X,B,H,P,M,C,R,S$ denoting context, bounded scope, hypothesis, pattern,
model implications, contrast assessment, refutation, and routing:
\[
\begin{aligned}
 e_0 &:X\to\mathsf{Gen}[B],\qquad
 a:B\to\mathsf{Gen}[H],\quad i:B\to\mathsf{Gen}[P],\\
 m &:J_{HP}\to\mathsf{Gen}[M],\qquad
 c:J_{HP}\to\mathsf{Gen}[C],\\
 r &:J_{MC}\to\mathsf{Gen}[R],\qquad
 e_1:J_{HPR}\to\mathsf{Gen}[S],\\
 J_{HP}&=\{\mathtt{discovery\_hypothesis}:H,\ \mathtt{discovery\_pattern}:P\},\\
 J_{MC}&=\{\mathtt{model\_implications}:M,\ \mathtt{contrast\_assessment}:C\},\\
 J_{HPR}&=\{\mathtt{discovery\_hypothesis}:H,\ \mathtt{discovery\_pattern}:P,
          \ \mathtt{refutation\_assessment}:R\}.
\end{aligned}
\]
Both formation calls consume $e_0(x).\mathtt{value}$. Both probes consume the
same named join of $a$'s and $i$'s projected values; $r$ receives the projected
probe outputs; $e_1$ retains the formation outputs alongside $r$'s value.
\emph{Scopes\allowbreak Discovery\allowbreak Formation}, \emph{Grounds\allowbreak Discovery\allowbreak Probe},
\emph{Grounds\allowbreak Discovery\allowbreak Refutation}, and \emph{Calibrates\allowbreak Discovery\allowbreak Routing}
respectively identify these four evidential dependencies, without establishing
that the model or comparison is adequate.

\paragraph{Constructed example.}
Within a fixed sensor dataset, induction finds a drifting offset and abduction
proposes either thermal sensitivity or clock drift. A model predicts offset
changes under temperature variation; the contrast holds temperature fixed while
varying elapsed time. If both hypotheses make the same prediction under available
controls, no run selects between them: the next task is a discriminating control.
If the offset changes only with elapsed time, thermal sensitivity is weakened
within that setup. The portfolio preserves the calibration target and records
which observation could reverse the decision.

The native graph has two parallel stages and returns Promotable, Reorient,
TargetReset, or Stalled. The original routes the first three to release,
conceptual exploration, or navigation, and stops on a stalled reasoning run.
It supplies the cross-cycle control policy; native output alternatives alone do
not execute a marathon or establish termination.

% SPDX-FileCopyrightText: © 2026 Idiap Research Institute <contact@idiap.ch>
%
% SPDX-FileContributor: Andre Freitas <andre.freitas@idiap.ch>
% Neuro-symbolic AI Group
%
% SPDX-FileContributor: Danilo Gusicuma <danilo.gusicuma@idiap.ch>
% Neuro-symbolic AI Group
%
% SPDX-License-Identifier: MIT
% Generated from step-explanations.json
\paragraph{Step-by-step reconstruction.}
A constructed sensor record contains temperature, elapsed time and a drifting offset. Thermal sensitivity and clock drift are rival explanations.
The following is a source-guided illustrative reconstruction. Internal procedures are prescriptions; the native declarations do not establish their execution.
\begin{enumerate}[leftmargin=*,label=\arabic*.]
\item \textbf{Bound the search} (\texttt{search\_\allowbreak{}scope}).
\begin{description}[style=standard,labelindent=0pt,leftmargin=5.5em,labelwidth=5em,labelsep=0.5em,align=left,font=\normalfont\bfseries,itemsep=0.15em,parsep=0pt,topsep=0.3em]
\item[Purpose] Bound admissible evidence and comparisons.
\item[Input] Sensor context; measurement constraints; calibration target
\item[Method] Separate available evidence from desired measurements and delimit the regime in which candidate explanations must operate.
\item[Check] A useful scope retains the target and exposes missing controls.
\item[Output] Bounded sensor record and admissible comparisons
\item[Next] Both formation branches consume this scope; insufficient controls constrain subsequent conclusions.
\item[Example] Freeze the measured offset, temperature and elapsed-time columns; note that no intervention has yet separated thermal and temporal effects.
\end{description}
\item \textbf{Propose a mechanism} (\texttt{discovery\_\allowbreak{}hypothesis}).
\begin{description}[style=standard,labelindent=0pt,leftmargin=5.5em,labelwidth=5em,labelsep=0.5em,align=left,font=\normalfont\bfseries,itemsep=0.15em,parsep=0pt,topsep=0.3em]
\item[Purpose] Propose mechanisms with distinguishable predictions.
\item[Input] Bounded sensor record, shared with pattern induction
\item[Method] Construct rival dependencies and specify the observable on which they could disagree, rather than naming a cause without a test.
\item[Check] A candidate must address the calibration target and admit a falsifier.
\item[Output] Candidate thermal mechanism; clock-drift rival
\item[Next] Both probes receive the hypothesis with the pattern; unsupported mechanisms remain candidates.
\item[Example] Thermal sensitivity predicts offset changes with temperature at fixed elapsed time; clock drift predicts changes with elapsed time at fixed temperature.
\end{description}
\item \textbf{Extract the observed regularity} (\texttt{discovery\_\allowbreak{}pattern}).
\begin{description}[style=standard,labelindent=0pt,leftmargin=5.5em,labelwidth=5em,labelsep=0.5em,align=left,font=\normalfont\bfseries,itemsep=0.15em,parsep=0pt,topsep=0.3em]
\item[Purpose] Describe observed variation without inferring causation.
\item[Input] Bounded observations; no hypothesis output
\item[Method] Compare offsets across the observed conditions and state the trend together with its observational limits.
\item[Check] The pattern must be supported within the record and retain confounding.
\item[Output] Provisional offset trend, with confounding retained
\item[Next] Join this pattern with the hypothesis for both probes; it is not independent confirmation.
\item[Example] The record shows increasing offset while both temperature and time increase. This trend fits both candidate mechanisms and selects neither.
\end{description}
\item \textbf{Derive model predictions} (\texttt{model\_\allowbreak{}implications}).
\begin{description}[style=standard,labelindent=0pt,leftmargin=5.5em,labelwidth=5em,labelsep=0.5em,align=left,font=\normalfont\bfseries,itemsep=0.15em,parsep=0pt,topsep=0.3em]
\item[Purpose] Derive comparable predictions under explicit assumptions.
\item[Input] Discovery hypothesis + provisional observed pattern
\item[Method] Specify variables, controls and measurement assumptions; preregister which outcomes favor, weaken or fail to distinguish each rival, including an inconclusive region, before deriving comparable predictions.
\item[Check] Predictions must concern the same observable and regime as the target.
\item[Output] Conditional thermal and temporal model predictions
\item[Next] Refutation later receives these implications with the contrast assessment; model failure remains informative.
\item[Example] Under thermal sensitivity, holding temperature fixed removes the predicted thermal offset change; under clock drift the offset can still grow with time.
\end{description}
\item \textbf{Design the discriminating contrast} (\texttt{contrast\_\allowbreak{}assessment}).
\begin{description}[style=standard,labelindent=0pt,leftmargin=5.5em,labelwidth=5em,labelsep=0.5em,align=left,font=\normalfont\bfseries,itemsep=0.15em,parsep=0pt,topsep=0.3em]
\item[Purpose] Identify a control that separates the rivals.
\item[Input] Discovery hypothesis + pattern, shared with model probe
\item[Method] Perturb one relevant assumption or control while preserving the others; explain which dependency the contrast tests.
\item[Check] A contrast is discriminating only if rivals predict different outcomes under justified controls.
\item[Output] Contrast assessment: informative control or non-identifiability
\item[Next] Send this assessment to refutation; unavailable controls require a new measurement plan.
\item[Example] Compare readings at fixed temperature but different elapsed times. Continued drift challenges a purely thermal account if other conditions remain stable.
\end{description}
\item \textbf{Combine the challenges} (\texttt{refutation\_\allowbreak{}assessment}).
\begin{description}[style=standard,labelindent=0pt,leftmargin=5.5em,labelwidth=5em,labelsep=0.5em,align=left,font=\normalfont\bfseries,itemsep=0.15em,parsep=0pt,topsep=0.3em]
\item[Purpose] Locate what survives both challenges.
\item[Input] Model implications + contrast assessment; scoped assumptions
\item[Method] Compare predicted differences with the available contrast and localize inconsistency; do not count two analyses of one record as two independent experiments.
\item[Check] A negative verdict needs a particular violated prediction, not an unfavorable label.
\item[Output] Refutation assessment: proposed falsifier, no observed refutation
\item[Next] Calibration receives this finding with the original hypothesis and pattern.
\item[Example] If a future fixed-temperature test finds drift, it would defeat the thermal-only prediction under its control assumptions. The current confounded record supplies no such finding.
\end{description}
\item \textbf{Allocate the next discovery action} (\texttt{routing\_\allowbreak{}assessment}).
\begin{description}[style=standard,labelindent=0pt,leftmargin=5.5em,labelwidth=5em,labelsep=0.5em,align=left,font=\normalfont\bfseries,itemsep=0.15em,parsep=0pt,topsep=0.3em]
\item[Purpose] Choose continuation, reorientation or reset.
\item[Input] Hypothesis + pattern + refutation; retained calibration target
\item[Method] Compare target relevance and what the next action could change; preserve alternatives that remain indistinguishable.
\item[Check] Continuation needs new learning or a sharper test, not repeated failure.
\item[Output] Routing assessment: remaining mechanism question and justification
\item[Next] Original cross-cycle routing is a policy; the native return does not execute a marathon.
\item[Example] Retain the trend and both rival explanations; request the fixed-temperature measurement. Demote thermal-only only if the predicted falsifier is subsequently observed under the stated controls.
\end{description}
\end{enumerate}
\paragraph{Changed-input contrast.}
If temperature and elapsed time are perfectly confounded, neither probe identifies the cause. The proper next action is a discriminating measurement, not promotion of the more elaborate explanation.


\subsection{Conceptual repair: turning an obstruction into a testable bridge}
\label{app:granular-conceptual-hypothesis-exploration-loop}
This workflow changes the representation used to attack a problem. Obstruction
assessment ($e_0$) identifies the structural pressure behind a failure. Mediation
formation ($a$) proposes what kind of concept could resolve it. Analogical
transfer ($t$) builds a correspondence rather than merely naming a similar
subject. Deduction ($d$) exposes a consequence of that bridge. Refutation ($r$)
tests the consequence, and routing ($e_1$) evaluates the bridge in light of that
test. The distinction between mediation and bridge prevents a promising analogy
from being treated as an already specified transport of reasoning.

\paragraph{What is revised, and why.}
The linked bridge-forging procedure searches for a missing carrier, invariant,
representation, interface, or transfer law where two structures nearly meet.
It must state why existing objects fail, how the new object would be constructed,
which invariants survive, and its first failure regime. Defining the desired
conclusion into the bridge is explicitly disallowed. The route ranker then
compares target contribution, validation strength, risk, and tractability; a
beautiful bridge without proof-form obligations stays outside the proof queue.
Its feedback decision distinguishes informative adaptation from backtracking,
pivoting, and quarantine, and requires the selected machinery to control the
actual bottleneck. The workflow contract retains failure cartography, mediation
candidates, and validation packets for this decision. This describes prescribed
conceptual work around the native ports, not a recovered implementation of their
string-bearing values.

Let $F,P,H,B,C,R,S$ name failure state, obstruction pressure, mediation
hypothesis, bridge, consequence, refutation, and refinement status.
\[
\begin{aligned}
 e_0 &:F\to\mathsf{Gen}[P],&a &:P\to\mathsf{Gen}[H],\\
 t &:H\to\mathsf{Gen}[B],&d &:B\to\mathsf{Gen}[C],\\
 r &:C\to\mathsf{Gen}[R],&e_1 &:J_{BR}\to\mathsf{Gen}[S],\\
 p&=e_0(f),\ h=a(p.\mathtt{value}),\ b=t(h.\mathtt{value}),\\
 q&=d(b.\mathtt{value}),\ v=r(q.\mathtt{value}),\\
 J_{BR}&=\{\mathtt{conceptual\_bridge}:B,\ \mathtt{refutation\_assessment}:R\},\\
 s&=e_1\bigl\{\mathtt{conceptual\_bridge}=b.\mathtt{value},\\
  &\hspace{3em}\mathtt{refutation\_assessment}=v.\mathtt{value}\bigr\}.
\end{aligned}
\]
The relation sequence bounds obstruction pressure, transfers mediation into a
bridge, treats the bridge as a consequence premise, tests that consequence,
and calibrates routing. In particular, the transfer relation records the proposed
correspondence; the assumption relation keeps its consequences conditional.
Neither declaration asserts that analogy preserves all relevant structure.

\paragraph{Constructed example.}
A convergence argument for a physical state fails because a norm counts changes
of gauge as changes of observable state. Obstruction assessment identifies that
identity mismatch. Mediation proposes a quotient by gauge equivalence; the bridge
must construct the projection and show that dynamics descend to the quotient.
Deduction can then formulate convergence of equivalence classes. If two equivalent
representatives evolve into inequivalent classes, refutation exposes failure of
well-defined dynamics. The next task is to repair that compatibility law, not to
rename the quotient as a solution. This constructed case illustrates how a
representation change generates a new, exact proof burden.

Native output distinguishes WorkReady, ReadyForReleasePromotion, and Refuted.
Original ERFL sends work-ready repairs back to abductive proof work, release
candidates to promotion, and stops on refutation. The native straight sequence
contains neither these switches nor the original early-stop guard. Its final
join retains the bridge and assessment; it has no separately typed retained port
for the derived consequence or original failure state.

% SPDX-FileCopyrightText: © 2026 Idiap Research Institute <contact@idiap.ch>
%
% SPDX-FileContributor: Andre Freitas <andre.freitas@idiap.ch>
% Neuro-symbolic AI Group
%
% SPDX-FileContributor: Danilo Gusicuma <danilo.gusicuma@idiap.ch>
% Neuro-symbolic AI Group
%
% SPDX-License-Identifier: MIT
% Generated from step-explanations.json
\paragraph{Step-by-step reconstruction.}
A state-convergence argument uses a norm that distinguishes gauge representatives even when they represent the same observable state.
The following is a source-guided illustrative reconstruction. Internal procedures are prescriptions; the native declarations do not establish their execution.
\begin{enumerate}[leftmargin=*,label=\arabic*.]
\item \textbf{Locate the representational obstruction} (\texttt{obstruction\_\allowbreak{}pressure}).
\begin{description}[style=standard,labelindent=0pt,leftmargin=5.5em,labelwidth=5em,labelsep=0.5em,align=left,font=\normalfont\bfseries,itemsep=0.15em,parsep=0pt,topsep=0.3em]
\item[Purpose] Distinguish conceptual obstruction from unfinished calculation.
\item[Input] Failed argument; norm’s identity distinction
\item[Method] Trace the failed implication to the role played by the carrier, invariant or equivalence relation.
\item[Check] An obstruction diagnosis must explain why existing objects fail at the named boundary.
\item[Output] Obstruction diagnosis: representative-dependent convergence criterion
\item[Next] Mediation consumes this diagnosis; a merely missing calculation should remain a derivation task.
\item[Example] Two gauge-equivalent representatives can be far apart in the chosen norm, so representative convergence need not express physical-state convergence.
\end{description}
\item \textbf{Propose the mediating object} (\texttt{mediation\_\allowbreak{}hypothesis}).
\begin{description}[style=standard,labelindent=0pt,leftmargin=5.5em,labelwidth=5em,labelsep=0.5em,align=left,font=\normalfont\bfseries,itemsep=0.15em,parsep=0pt,topsep=0.3em]
\item[Purpose] Propose a carrier for the intended question.
\item[Input] Representative-dependent identity and convergence obstruction
\item[Method] Propose a quotient or other mediator, naming what it would preserve and its first possible failure regime.
\item[Check] The desired convergence theorem cannot be hidden inside the mediator's definition.
\item[Output] Mediation hypothesis: gauge-equivalence classes
\item[Next] Transfer must turn this proposal into specified maps; a name alone is insufficient.
\item[Example] Replace the proposed carrier of raw representatives by classes under gauge equivalence, while leaving existence of suitable dynamics as an obligation.
\end{description}
\item \textbf{Specify the bridge} (\texttt{conceptual\_\allowbreak{}bridge}).
\begin{description}[style=standard,labelindent=0pt,leftmargin=5.5em,labelwidth=5em,labelsep=0.5em,align=left,font=\normalfont\bfseries,itemsep=0.15em,parsep=0pt,topsep=0.3em]
\item[Purpose] Specify the mediator’s maps and compatibility burden.
\item[Input] Quotient proposal; intended observable interpretation
\item[Method] Specify projection from representatives to classes and the compatibility needed for dynamics to descend.
\item[Check] Equivalent representatives must induce the same class evolution for the bridge to be well defined.
\item[Output] Projection bridge; open dynamics-compatibility obligation
\item[Next] Deduction uses this bridge conditionally; an absent compatibility law blocks unconditional transport.
\item[Example] For representative evolution F, require the class of F(x) to equal the class of F(y) whenever x and y are gauge equivalent.
\end{description}
\item \textbf{Derive a bridge consequence} (\texttt{bridge\_\allowbreak{}consequence}).
\begin{description}[style=standard,labelindent=0pt,leftmargin=5.5em,labelwidth=5em,labelsep=0.5em,align=left,font=\normalfont\bfseries,itemsep=0.15em,parsep=0pt,topsep=0.3em]
\item[Purpose] State the bridge’s conditional consequence.
\item[Input] Mediating bridge; open compatibility assumptions
\item[Method] Formulate the transported statement on the new carrier and retain each bridge premise.
\item[Check] The consequence must concern equivalence classes without silently restoring representative identity.
\item[Output] Class-level evolution or convergence claim, conditional on its premises
\item[Next] Refutation tests this consequence; convergence estimates remain separate from well-definedness.
\item[Example] A quotient trajectory can be defined from F only if changing the representative leaves its next class unchanged; convergence would then be a class-level question.
\end{description}
\item \textbf{Test preservation} (\texttt{refutation\_\allowbreak{}assessment}).
\begin{description}[style=standard,labelindent=0pt,leftmargin=5.5em,labelwidth=5em,labelsep=0.5em,align=left,font=\normalfont\bfseries,itemsep=0.15em,parsep=0pt,topsep=0.3em]
\item[Purpose] Test whether the representation is well defined.
\item[Input] Bridge consequence; compatibility premise
\item[Method] Search for equivalent inputs whose evolved outputs cease to be equivalent; localize the precise preservation law at fault.
\item[Check] A bridge fails this test when one physical input class receives two distinct output classes.
\item[Output] Refutation assessment: well-definedness defect or bounded surviving claim
\item[Next] Pass the finding to routing with the retained bridge; no found defect is not a proof.
\item[Example] Choose equivalent x and y with inequivalent F(x) and F(y). The proposed class dynamics is ambiguous, independently of any convergence estimate.
\end{description}
\item \textbf{Route the conceptual repair} (\texttt{refinement\_\allowbreak{}routing}).
\begin{description}[style=standard,labelindent=0pt,leftmargin=5.5em,labelwidth=5em,labelsep=0.5em,align=left,font=\normalfont\bfseries,itemsep=0.15em,parsep=0pt,topsep=0.3em]
\item[Purpose] Choose proof work or bridge reconstruction.
\item[Input] Conceptual bridge + its refutation assessment
\item[Method] Compare target contribution, validation strength and the exact preservation defect; distinguish an informative repair from a repeated blocked route.
\item[Check] A work-ready bridge needs explicit obligations that address the bottleneck.
\item[Output] Refinement assessment: compatibility-repair obligation
\item[Next] Original policy routes work or release and stops refutation; native routing is a returned value.
\item[Example] Keep quotienting as a candidate representation but require a compatible evolution law; do not advertise the quotient name as a solved convergence problem.
\end{description}
\end{enumerate}
\paragraph{Changed-input contrast.}
If equivalent representatives evolve into equivalent classes, the proposed quotient passes this well-definedness test. Convergence still requires its own argument; compatibility alone does not establish it.


\subsection{Frontier deepening: selecting before formalizing}
\label{app:granular-conceptual-frontier-deepening}
A concept should become a formal claim only after its discriminating content has
survived an explicit selection step. Frontier grounding ($e_0$) fixes a bounded
search from frozen context. Hypothesis formation ($a$) proposes a candidate
within it. Refutation ($r$) challenges that candidate before selection ($e_1$)
compares its promise with the challenge. Bridge revision ($t$) turns the selected
content into a correspondence usable for formal work; deduction ($d$) expresses
a formal consequence. Readiness calibration ($e_2$) decides whether that claim
is ready for theorem work, not whether it is proved.

\paragraph{The conceptual work hidden by a single hypothesis port.}
The linked frontier procedure opens invention before assigning family labels:
authors work without sibling candidates, then a separate assessor reconstructs
families from primitive changes, mechanism, intervention, prediction, and first
falsifier. Semantic paraphrases collapse; incompatible mechanisms remain live.
Reflection changes identity conditions, exposes internal contradictions and
inherited commitments, separates necessary from regime-relative structure, and
seeks unification that preserves decisive differences. Blind judgment compares
complete candidates on non-compensatory dimensions such as relevance,
protected-success preservation, explanatory depth, correction, and falsifiability;
a plural non-dominated frontier is an admissible outcome. Category or HoTT
representations earn a role only by exposing an actual mediator or transport
obstruction. This prescribed process is much richer than the linear native
projection. There is also a source difference: the skill makes formal routes
conditional, while the named contract lists central-route derivation and formal
claim production as required steps. Neither should be silently read as proof
that every candidate becomes formal-ready.

The local aliases are $X$ (frozen context), $F$ (frontier), $H$ (hypothesis),
$R$ (refutation), $S$ (selection), $B$ (revised bridge), $C$ (formal claim), and
$Q$ (readiness).
\[
\begin{aligned}
 e_0 &:X\to\mathsf{Gen}[F],&a &:F\to\mathsf{Gen}[H],\\
 r &:H\to\mathsf{Gen}[R],&e_1 &:J_{HR}\to\mathsf{Gen}[S],\\
 t &:S\to\mathsf{Gen}[B],&d &:B\to\mathsf{Gen}[C],\\
 e_2 &:C\to\mathsf{Gen}[Q],\\
 J_{HR}&=\{\mathtt{conceptual\_hypothesis}:H,
           \ \mathtt{refutation\_assessment}:R\}.
\end{aligned}
\]
Every sequential call uses the predecessor's $\mathtt{.value}$; selection instead
joins the retained hypothesis value with its refutation value. The search-bound
relation attaches formation to the grounded frontier. The test relation attaches
criticism to the candidate, and the calibration relation is used twice: first
for selection, then for readiness. Bridge revision and premise relations mark
the remaining changes of commitment. Reusing a relation name does not equate
these two calibration judgments.

\paragraph{Constructed example.}
For two effective descriptions with different state variables but matching
measured predictions in a bounded regime, invention may propose a common latent
carrier or a direct translation between observables. Reflection asks which
identity distinctions each proposal erases. A candidate transport must specify
maps, the regime in which they apply, and a test where predictions diverge.
Agreement on existing observations does not establish equivalence, and a latent
carrier without a discriminating consequence receives no automatic preference.
Both candidates can remain on the frontier until a correction or new observable
distinguishes them. The formalization task, when justified, is the specific
transport or compatibility obligation.

The native output offers FormalReady, NotFormalReadyOrNonDominated, or Refuted,
but the native sequence has no switch. Original ERFL routes readiness to theorem
proof, returns a non-ready or non-dominated candidate to exploration, and stops
at a refuted $r$. Moreover, selection, bridge, and claim are nominal string-bearing
types: their signatures do not enforce preservation of the original evidence or
hypothesis identity. Such stronger indexed obligations must be proposed
explicitly rather than attributed to this native declaration.

% SPDX-FileCopyrightText: © 2026 Idiap Research Institute <contact@idiap.ch>
%
% SPDX-FileContributor: Andre Freitas <andre.freitas@idiap.ch>
% Neuro-symbolic AI Group
%
% SPDX-FileContributor: Danilo Gusicuma <danilo.gusicuma@idiap.ch>
% Neuro-symbolic AI Group
%
% SPDX-License-Identifier: MIT
% Generated from step-explanations.json
\paragraph{Step-by-step reconstruction.}
Two effective descriptions match existing measurements in a bounded regime but use different state variables. A latent carrier and a direct observable translation are candidate bridges.
The following is a source-guided illustrative reconstruction. Internal procedures are prescriptions; the native declarations do not establish their execution.
\begin{enumerate}[leftmargin=*,label=\arabic*.]
\item \textbf{Freeze the conceptual pressure} (\texttt{conceptual\_\allowbreak{}frontier}).
\begin{description}[style=standard,labelindent=0pt,leftmargin=5.5em,labelwidth=5em,labelsep=0.5em,align=left,font=\normalfont\bfseries,itemsep=0.15em,parsep=0pt,topsep=0.3em]
\item[Purpose] Fix protected successes and the unresolved distinction.
\item[Input] Frozen descriptions; bounded empirical agreement
\item[Method] Separate target commitments from replaceable concepts and identify what a new bridge must explain.
\item[Check] The frontier must retain regime limits and the successes a replacement may not discard.
\item[Output] Bounded frontier: relation between descriptions; protected distinctions
\item[Next] Hypothesis formation consumes this frontier; ordinary missing calculation does not justify a new ontology.
\item[Example] Record agreement on current observables, different state variables, and an open question about predictions beyond that observable set.
\end{description}
\item \textbf{Invent alternatives} (\texttt{conceptual\_\allowbreak{}hypothesis}).
\begin{description}[style=standard,labelindent=0pt,leftmargin=5.5em,labelwidth=5em,labelsep=0.5em,align=left,font=\normalfont\bfseries,itemsep=0.15em,parsep=0pt,topsep=0.3em]
\item[Purpose] Generate meaningfully different explanatory structures.
\item[Input] Bounded frontier; protected commitments
\item[Method] The linked skill prescribes 8--12 isolated raw candidates, blinded reconstruction of at least six semantic families unless verifier-replayed constraint evidence justifies fewer, and six reflective transformations. Propose primitive, mediator or transport changes and candidate-specific falsifiers; this two-candidate example illustrates one distinction only.
\item[Check] Semantic paraphrases do not create new families; each alternative must change a meaningful dependency.
\item[Output] Candidate explanations: latent mediation or direct observable translation
\item[Next] Refutation examines candidate content; a nominal hypothesis port does not implement the larger portfolio procedure.
\item[Example] A latent carrier attempts to explain both state descriptions, while direct translation relates selected observables without asserting a common underlying state space.
\end{description}
\item \textbf{Challenge a candidate's distinctions} (\texttt{refutation\_\allowbreak{}assessment}).
\begin{description}[style=standard,labelindent=0pt,leftmargin=5.5em,labelwidth=5em,labelsep=0.5em,align=left,font=\normalfont\bfseries,itemsep=0.15em,parsep=0pt,topsep=0.3em]
\item[Purpose] Test what the proposed bridge fails to preserve.
\item[Input] Candidate explanation; claimed regime
\item[Method] Identify a decisive observable, intervention or identity distinction and ask whether the proposal preserves it.
\item[Check] A criticism must change a dependency or reveal an unmet prediction, not merely request different wording.
\item[Output] Refutation finding: unresolved discrimination or specific candidate defect
\item[Next] Selection receives the original hypothesis with this finding, preserving its identity.
\item[Example] Ask whether two states merged by the latent carrier could disagree on a newly measured observable. If so, that merging needs justification or repair.
\end{description}
\item \textbf{Select without forcing consensus} (\texttt{hypothesis\_\allowbreak{}selection}).
\begin{description}[style=standard,labelindent=0pt,leftmargin=5.5em,labelwidth=5em,labelsep=0.5em,align=left,font=\normalfont\bfseries,itemsep=0.15em,parsep=0pt,topsep=0.3em]
\item[Purpose] Select defensible content without forcing preference.
\item[Input] Hypothesis + its own refutation assessment
\item[Method] The linked skill prescribes preregistered non-compensatory dimensions, at least three separately configured blinded judges and sham/rival controls. Compare complete candidates and retain a nondominated set; missing judges, warrants or failed controls yield a typed null rather than a forced preference.
\item[Check] Failure on a required dimension cannot be offset by elegance elsewhere.
\item[Output] Selected content; unresolved alternatives where warranted
\item[Next] Bridge revision consumes selected content; non-ready plural results have an original exploration route.
\item[Example] Keep both latent mediation and direct translation if current evidence cannot distinguish them; make the missing observable test part of their retained burden.
\end{description}
\item \textbf{Make a selected bridge explicit} (\texttt{conceptual\_\allowbreak{}bridge}).
\begin{description}[style=standard,labelindent=0pt,leftmargin=5.5em,labelwidth=5em,labelsep=0.5em,align=left,font=\normalfont\bfseries,itemsep=0.15em,parsep=0pt,topsep=0.3em]
\item[Purpose] Specify a useful formal correspondence.
\item[Input] Selected content; open identity and regime obligations
\item[Method] Specify maps, carriers and preservation conditions only where formal representation adds a concrete obligation.
\item[Check] Category or type language must expose a mediator or transport obstruction, not supply prestige.
\item[Output] Revised bridge: specified compatibility obligation
\item[Next] Deduction uses this bridge; linked procedures make formal routes conditional although the native sequence is linear.
\item[Example] For the direct-translation candidate, specify an observable map and demand agreement of predictions over the chosen regime, without assuming full state-space equivalence.
\end{description}
\item \textbf{State the exact formal burden} (\texttt{formal\_\allowbreak{}claim}).
\begin{description}[style=standard,labelindent=0pt,leftmargin=5.5em,labelwidth=5em,labelsep=0.5em,align=left,font=\normalfont\bfseries,itemsep=0.15em,parsep=0pt,topsep=0.3em]
\item[Purpose] Formulate the bridge’s precise proof obligation.
\item[Input] Revised map; preservation conditions
\item[Method] Express the compatibility statement with explicit domain and quantifiers, retaining assumptions as premises.
\item[Check] The claim must not strengthen regime-limited prediction agreement to unrestricted equivalence.
\item[Output] Formal map-compatibility claim; retained premises and regime
\item[Next] Readiness assesses this statement; the act of stating it does not discharge it.
\item[Example] The obligation becomes prediction preservation for the specified observable map on the bounded regime; unseen observables remain outside the claim.
\end{description}
\item \textbf{Assess readiness for theorem work} (\texttt{calibrated\_\allowbreak{}candidate}).
\begin{description}[style=standard,labelindent=0pt,leftmargin=5.5em,labelwidth=5em,labelsep=0.5em,align=left,font=\normalfont\bfseries,itemsep=0.15em,parsep=0pt,topsep=0.3em]
\item[Purpose] Assess proof readiness while retaining unresolved choices.
\item[Input] Formal claim; domain; open obligations
\item[Method] Check that objects and the required implication are definite; distinguish readiness from an established theorem.
\item[Check] An unresolved choice of carrier or target statement prevents an unconditional readiness judgment.
\item[Output] Calibrated candidate: ready, non-ready or unresolved plural status
\item[Next] Original policy can send ready claims to theorem work; the native sequence supplies no switch.
\item[Example] If the observable map and its regime are specified, queue its compatibility lemma; if the carrier remains ambiguous, return to conceptual selection.
\end{description}
\end{enumerate}
\paragraph{Changed-input contrast.}
If a new observable distinguishes the candidates, selection must incorporate that difference. Agreement on the original measurements alone warrants neither equivalence nor a unique winner.


\subsection{Evidence-to-proof onramp: finding the obligation worth proving}
\label{app:granular-evidence-to-proof-route-onramp}
The onramp separates proposing a useful route from selecting an exact entrypoint.
Grounding ($e_0$) anchors a source bundle in evidence. Route abduction ($a_0$)
proposes how that evidence could support proof work. Deduction ($d_0$) makes a
consequence inspectable, and calibration ($e_1$) assesses it. A second abductive
occurrence ($a_1$) then revises the proof-obligation entrypoint: assessment may
show that a different lemma, rather than the original ambitious target, is the
tractable next step. Refutation ($r_0$) challenges that entrypoint; final
calibration ($e_2$) preserves it with its falsification record.

\paragraph{How an entrypoint is found.}
The named contract orders ingestion and theorem normalization before conjecture
grounding and connection search. The connector distinguishes established
statement dependencies from proposed cross-component links, using symbol roles,
method overlap, and candidate reductions to expose missing bridges. The hub audit
asks whether a highly connected definition is being mistaken for mathematical
progress. Pathfinding then anchors each agenda on one statement, attaches its
supporting hypotheses and links, and gives it a stop or pruning condition. A
regime-crossing route additionally needs an audit of the transfer; the final
handoff retains the route together with its falsification evidence. Consequently,
route selection concerns the missing implication that would advance the target,
not simply the most connected node. These are ordered contract and prescribed
skill procedures; the seven native occurrences are their coarse conceptual
projection, not an execution trace of all those steps.

For source $X$, ground $G$, route $H$, consequence $C$, assessment $A$, entrypoint
$P$, falsification $R$, and readiness $S$:
\[
\begin{aligned}
 e_0 &:X\to\mathsf{Gen}[G],&a_0 &:G\to\mathsf{Gen}[H],\\
 d_0 &:H\to\mathsf{Gen}[C],&e_1 &:C\to\mathsf{Gen}[A],\\
 a_1 &:A\to\mathsf{Gen}[P],&r_0 &:P\to\mathsf{Gen}[R],\\
 e_2 &:J_{PR}\to\mathsf{Gen}[S],\\
 J_{PR}&=\{\mathtt{proof\_entrypoint}:P,
           \ \mathtt{falsification\_assessment}:R\}.
\end{aligned}
\]
After grounding, successive unary calls pass explicit $\mathtt{.value}$
projections along the sequence; $e_2$ joins the projected output of $a_1$ with
the projected output of $r_0$ in the named fields.
Relations identify grounding, conditional consequence, calibration,
entrypoint revision, and testing; calibration is reused at both assessments.
The readiness record contains a disposition, expert dossier, falsification
assessment, and hub-bias audit. The audit is an output field, not an additional
declared reasoning occurrence.

\paragraph{Constructed example.}
A knowledge base connects a circuit lower-bound target to a highly central
complexity-class definition and to a less connected lemma translating proofs into
circuits. The hub audit prevents the definition's degree from deciding the route.
Pathfinding instead anchors on the transfer lemma and asks whether it preserves
uniformity. A counterexample search tests whether the translation relies on
nonuniform advice; if it does, that route cannot support the uniform target
without a new lemma. The dossier therefore carries an explicit uniformity
obligation or a stopped route, rather than a graph-centrality score presented as
proof progress. This is a constructed route diagnosis, not a lower-bound result.

Original ERFL switches a ready route to theorem delivery, sends conceptual gaps
to exploration, and stops on refutation. Native comments describe these
handoffs, including an adapter to \texttt{ProofQueueItem}, as non-executable.
The native sequence itself returns the assessment; it neither implements that
adapter nor enforces the early stop. The nominal ID-bearing intermediate records
also leave evidence-preservation obligations outside their signatures.

% SPDX-FileCopyrightText: © 2026 Idiap Research Institute <contact@idiap.ch>
%
% SPDX-FileContributor: Andre Freitas <andre.freitas@idiap.ch>
% Neuro-symbolic AI Group
%
% SPDX-FileContributor: Danilo Gusicuma <danilo.gusicuma@idiap.ch>
% Neuro-symbolic AI Group
%
% SPDX-License-Identifier: MIT
% Generated from step-explanations.json
\paragraph{Step-by-step reconstruction.}
A lower-bound knowledge base contains a highly connected complexity-class definition and a less connected proposed proof-to-circuit translation. The target requires uniformity.
The following is a source-guided illustrative reconstruction. Internal procedures are prescriptions; the native declarations do not establish their execution.
\begin{enumerate}[leftmargin=*,label=\arabic*.]
\item \textbf{Ground the route evidence} (\texttt{e0}).
\begin{description}[style=standard,labelindent=0pt,leftmargin=5.5em,labelwidth=5em,labelsep=0.5em,align=left,font=\normalfont\bfseries,itemsep=0.15em,parsep=0pt,topsep=0.3em]
\item[Purpose] Ground statements before proposing proof connections.
\item[Input] Source bundle; target; source references
\item[Method] Normalize theorem bodies and retain hypotheses, provenance and the status of each imported claim.
\item[Check] Only licensed statements may serve as established premises; topic proximity is insufficient.
\item[Output] Grounded definitions; proposed translation claims with status
\item[Next] Route formation consumes this evidence; ungrounded statements remain conjectural or excluded.
\item[Example] Record the uniform lower-bound target, the central definition, and a translation candidate whose uniformity condition has not been checked.
\end{description}
\item \textbf{Propose a useful connection} (\texttt{a0}).
\begin{description}[style=standard,labelindent=0pt,leftmargin=5.5em,labelwidth=5em,labelsep=0.5em,align=left,font=\normalfont\bfseries,itemsep=0.15em,parsep=0pt,topsep=0.3em]
\item[Purpose] Identify a target-relevant missing implication.
\item[Input] Grounded statements; dependency context
\item[Method] Compare candidate reductions and cross-component links using symbol roles and method overlap, separating proposed edges from established dependencies.
\item[Check] A route must supply proof relevance beyond graph centrality.
\item[Output] Route hypothesis: candidate proof-to-circuit translation
\item[Next] Deduction exposes its consequence; the selected bridge is still hypothetical.
\item[Example] Choose the less central translation because, if uniform, it could connect a proof lower bound to the circuit statement required by the target.
\end{description}
\item \textbf{Expose the route's consequence} (\texttt{d0}).
\begin{description}[style=standard,labelindent=0pt,leftmargin=5.5em,labelwidth=5em,labelsep=0.5em,align=left,font=\normalfont\bfseries,itemsep=0.15em,parsep=0pt,topsep=0.3em]
\item[Purpose] Expose the candidate bridge’s mathematical payoff.
\item[Input] Translation hypothesis; retained conditions
\item[Method] Derive the proposed implication while tracking the target quantifier that the translation must preserve.
\item[Check] Do not replace a uniform target with a nonuniform consequence.
\item[Output] Translation consequence, conditional on required uniformity
\item[Next] Assessment checks this consequence's relevance and missing support.
\item[Example] The route would transfer the chosen lower-bound obligation only if one uniform construction supplies the translated objects across input sizes.
\end{description}
\item \textbf{Assess the consequence} (\texttt{e1}).
\begin{description}[style=standard,labelindent=0pt,leftmargin=5.5em,labelwidth=5em,labelsep=0.5em,align=left,font=\normalfont\bfseries,itemsep=0.15em,parsep=0pt,topsep=0.3em]
\item[Purpose] Identify the decisive gap in the payoff.
\item[Input] Derived consequence; uniformity assumption
\item[Method] Compare the consequence to the original target and inspect whether graph prominence concealed a missing theorem.
\item[Check] A highly connected definition does not compensate for an unproved translation property.
\item[Output] Consequence assessment: missing uniformity-preservation support
\item[Next] Entrypoint revision consumes this diagnosis instead of repeating the full ambitious target.
\item[Example] The definition contributes vocabulary, but the route's progress depends on proving that translation does not use size-specific advice.
\end{description}
\item \textbf{Revise the entrypoint} (\texttt{a1}).
\begin{description}[style=standard,labelindent=0pt,leftmargin=5.5em,labelwidth=5em,labelsep=0.5em,align=left,font=\normalfont\bfseries,itemsep=0.15em,parsep=0pt,topsep=0.3em]
\item[Purpose] Reduce the route to a testable obligation.
\item[Input] Consequence assessment; localized uniformity gap
\item[Method] Anchor an agenda on one statement, attach its hypotheses and links, and specify a stop or pruning condition.
\item[Check] The entrypoint must directly control the missing implication.
\item[Output] Entrypoint hypothesis: uniformity-preservation lemma
\item[Next] Refutation challenges this entrypoint before a proof handoff is justified.
\item[Example] Replace the full lower-bound demand with the lemma that the proposed translation uses no nonuniform advice as input size varies.
\end{description}
\item \textbf{Challenge the entrypoint} (\texttt{r0}).
\begin{description}[style=standard,labelindent=0pt,leftmargin=5.5em,labelwidth=5em,labelsep=0.5em,align=left,font=\normalfont\bfseries,itemsep=0.15em,parsep=0pt,topsep=0.3em]
\item[Purpose] Test the selected lemma in its stated regime.
\item[Input] Uniformity-preservation entrypoint; hypotheses
\item[Method] Search for a translation instance requiring advice or a regime change; distinguish an actual counterexample from a missing argument.
\item[Check] A falsifier must attack the selected lemma, not a different complexity claim.
\item[Output] Falsification finding: advice dependence or unresolved uniformity proof
\item[Next] Final readiness joins this finding with the same entrypoint.
\item[Example] A translation that chooses unrelated advice at each input size defeats the proposed uniformity lemma, leaving any nonuniform statement separate.
\end{description}
\item \textbf{Bound the proof handoff} (\texttt{e2}).
\begin{description}[style=standard,labelindent=0pt,leftmargin=5.5em,labelwidth=5em,labelsep=0.5em,align=left,font=\normalfont\bfseries,itemsep=0.15em,parsep=0pt,topsep=0.3em]
\item[Purpose] Assess a proof task with its criticism intact.
\item[Input] Entrypoint hypothesis + its falsification assessment
\item[Method] Retain the route dossier, hub-bias diagnosis and transfer burden when judging readiness.
\item[Check] Readiness cannot arise solely from connectivity or a favorable label.
\item[Output] Readiness assessment: bounded proof task or reorientation rationale
\item[Next] Original routing needs a separate ProofQueueItem adapter; the native return does not construct it.
\item[Example] Hand off a precise unresolved uniformity obligation if viable, or stop the advice-dependent route without claiming a lower bound.
\end{description}
\end{enumerate}
\paragraph{Changed-input contrast.}
If the translation is independently shown to preserve uniformity, it can support a suitable proof entrypoint. Merely increasing the centrality of the definition never supplies that implication.


\subsection{Theorem delivery: preserving the exact claim through challenge}
\label{app:granular-theorem-proof-delivery}
The theorem workflow begins after a proof obligation has been chosen. Deduction
($d_0$) constructs an exact theorem claim from explicit premises. Refutation
($r_0$) tests that claim rather than a weaker paraphrase. Calibration ($e_0$)
combines the retained claim and refutation assessment into a disposition with a
formal-proof gate and, where applicable, a failure certificate. This is a
three-occurrence closure attempt: choosing a new hypothesis is deliberately not
a local occurrence. A failed attempt can instead supply conceptual work with an
exact statement and its unresolved obligations.

\paragraph{What happens inside a proof attempt.}
The linked builder orders lemmas without circular dependencies and records where
each premise enters. A parameter ledger propagates constants, exponents,
quantifiers, norms, and validity regimes so that a locally plausible step cannot
silently change the theorem. When a lemma blocks, the conditional gap-splitting
procedure identifies whether the cause is a missing estimate, transfer, wrong
quantifier, model mismatch, or false statement. It actively searches small,
degenerate, proxy, and adversarial cases before classifying the problem as missing
technique, then permits one explicit repair before choosing a new disposition.
The gate distinguishes an unconditional proof from an explicitly conditional
theorem and rejects silent gap closure. The contract therefore retains obligations
even when exit is allowed. These prescribed proof procedures supply substance to
``derive, test, calibrate'' without asserting that the three native operations
execute them or that a status value constitutes a formal proof certificate.

Write $P,C,R,S$ for explicit premises, exact claim, refutation, and disposition.
The input $q$ has native type \texttt{ProofQueueItem}.
\[
\begin{aligned}
 d_0 &:P\to\mathsf{Gen}[C],\qquad
 r_0:C\to\mathsf{Gen}[R],\qquad e_0:J_{CR}\to\mathsf{Gen}[S],\\
 c&=d_0(q.\mathtt{premises}),\qquad v=r_0(c.\mathtt{value}),\\
 J_{CR}&=\{\mathtt{theorem\_claim}:C,
           \ \mathtt{refutation\_assessment}:R\},\\
 s&=e_0\{\mathtt{theorem\_claim}=c.\mathtt{value},\
         \mathtt{refutation\_assessment}=v.\mathtt{value}\}.
\end{aligned}
\]
\emph{Tests\allowbreak Exact\allowbreak Theorem\allowbreak Claim} connects the challenge to the generated claim;
\emph{Calibrates\allowbreak Theorem\allowbreak Proof\allowbreak Disposition} connects the joint evidence to status.
The claim record explicitly retains premises and formal obligations. The output
retains claim, refutation, gate, and an optional failure certificate. These are
useful dependencies, but their nominal types do not prove that the gate is sound,
that obligations are discharged, or that the returned statement equals the
requested statement. A dependent specification could index those obligations by
$q$; this diagram reports the existing signature.

\paragraph{Constructed example.}
The queue requests exponential decay for a finite connected undirected graph's symmetric Laplacian on
the mean-zero subspace for $\dot u=-Lu$, assuming a gap
$\lambda_1\geq c>0$. Deduction retains these premises in the Euclidean-norm
claim $\|u(t)\|_2\leq e^{-ct}\|u(0)\|_2$ for $t\geq0$.
Refutation checks the excluded constant mode and the gap assumption. A missing
spectral argument remains an open obligation; favorable testing cannot close it.
The gate must distinguish an accepted proof from a certificate explaining that
gap or quarantine.

Original ERFL routes accepted proof status to contribution cultivation, failure
certificates to conceptual exploration, and quarantine to stop. Native comments
preserve that policy as non-executable handoffs, including adaptation to a result
atlas. The native sequence supplies no local terminal acceptance transition.

% SPDX-FileCopyrightText: © 2026 Idiap Research Institute <contact@idiap.ch>
%
% SPDX-FileContributor: Andre Freitas <andre.freitas@idiap.ch>
% Neuro-symbolic AI Group
%
% SPDX-FileContributor: Danilo Gusicuma <danilo.gusicuma@idiap.ch>
% Neuro-symbolic AI Group
%
% SPDX-License-Identifier: MIT
% Generated from step-explanations.json
\paragraph{Step-by-step reconstruction.}
Prove Euclidean-norm decay for the symmetric combinatorial Laplacian of a finite connected undirected graph on mean-zero initial data, assuming its first positive eigenvalue is at least c greater than zero.
The following is a source-guided illustrative reconstruction. Internal procedures are prescriptions; the native declarations do not establish their execution.
\begin{enumerate}[leftmargin=*,label=\arabic*.]
\item \textbf{Construct the exact conditional theorem} (\texttt{d0}).
\begin{description}[style=standard,labelindent=0pt,leftmargin=5.5em,labelwidth=5em,labelsep=0.5em,align=left,font=\normalfont\bfseries,itemsep=0.15em,parsep=0pt,topsep=0.3em]
\item[Purpose] Derive the exact queued obligation.
\item[Input] item.premises: queued exact\_\allowbreak{}statement; connected-graph, subspace and gap premises
\item[Method] Order lemmas without circular dependence and propagate constants, norms and quantifiers through each inference.
\item[Check] Every inferential step must identify its premises and preserve the theorem's scope.
\item[Output] Exact claim; retained premises and undischarged formal obligations
\item[Next] Refutation receives this claim; a missing lemma remains visible rather than silently closed.
\item[Example] For du/dt=-Lu on the connected graph, expand the mean-zero vector in positive-eigenvalue modes. Each coefficient evolves by exp(-lambda t), giving Euclidean-norm decay at rate at least c.
\end{description}
\item \textbf{Search for a defeater} (\texttt{r0}).
\begin{description}[style=standard,labelindent=0pt,leftmargin=5.5em,labelwidth=5em,labelsep=0.5em,align=left,font=\normalfont\bfseries,itemsep=0.15em,parsep=0pt,topsep=0.3em]
\item[Purpose] Challenge edge cases without weakening the claim.
\item[Input] Exact claim; gap premise; open obligations
\item[Method] Test small, degenerate and adversarial cases; classify a failed step as a false statement, wrong quantifier, model mismatch or missing technique.
\item[Check] Favorable examples cannot substitute for a missing spectral argument.
\item[Output] Refutation assessment: scoped surviving theorem or localized defect
\item[Next] Calibration retains this finding with the exact claim; repair must name any changed premise.
\item[Example] The constant mode violates unrestricted decay but is excluded by the mean-zero premise. Without the gap premise, cycle families defeat a common positive rate.
\end{description}
\item \textbf{Calibrate the proof disposition} (\texttt{e0}).
\begin{description}[style=standard,labelindent=0pt,leftmargin=5.5em,labelwidth=5em,labelsep=0.5em,align=left,font=\normalfont\bfseries,itemsep=0.15em,parsep=0pt,topsep=0.3em]
\item[Purpose] Distinguish proof discharge from conditional or blocked work.
\item[Input] Exact theorem claim + its refutation assessment
\item[Method] Check the lemma chain and retained obligations; issue a bounded disposition and failure certificate when a bridge remains missing.
\item[Check] A proof gate cannot discharge an obligation merely because refutation found no counterexample.
\item[Output] Claim; finding; proof gate; optional failure certificate
\item[Next] Original policy routes accepted work or failure; the native return is no terminal proof acceptance act.
\item[Example] Accept only the explicitly justified conditional spectral estimate; if spectral decomposition has not been supplied, retain that proof debt in the disposition.
\end{description}
\end{enumerate}
\paragraph{Changed-input contrast.}
If the queue requests all initial data, a nonzero constant vector defeats positive-rate decay. The theorem must retain the mean-zero restriction; an assessment of that restricted theorem cannot certify the broader one.

